\section{The catalog}
\label{sec:catalog}

Section~\ref{sec:field} asked a layer for two things: a phase function $\ph$ whose
level sets are the members, and its gradient magnitude, which converts a phase
residual into a distance. This section is the list, and it is deliberately short:
thirteen families, of which nine are one line of algebra each. The four that
are not---the radial pencil and the three lattices---and the reason they
are not, set up the rest of the paper.
The supplemental material (\S S1) draws all thirteen, each alone and then beating
against a perturbed copy of itself, which is the fastest way to see what the
algebra below buys.

\subsection{Nine closed forms}

For a family whose members are $\{\ph=ns+\varphi\}$, the index and distance fields
of \S\ref{sec:field} are
\begin{equation}
  \idx(p)=\frac{\ph(p)-\varphi}{s},
  \qquad
  d(p)=\frac{\mathrm{periodicDist}\!\left(\ph(p)-\varphi,\,s\right)}{\lvert\nabla\ph(p)\rvert},
  \label{eq:closedpair}
\end{equation}
so a family is completely specified by the pair $(\ph,\lvert\nabla\ph\rvert)$.
Table~\ref{tab:catalog} gives both for all nine (ten rows: three are written-out instances of the $N$-gon). Every entry is $O(1)$, branch-free
apart from the free-side $N$-gon, and calls no transcendental function except where
the geometry genuinely contains an angle.

\begin{table*}[t]
  \centering
  \caption{\figlead{The nine phase families.} Each row is a complete
  specification: substitute $\ph$ and $\lvert\nabla\ph\rvert$ into
  \eqref{eq:closedpair} and the layer renders; the square, triangle and
  hexagon rows are the $N=4$, $3$, $6$ instances of the $N$-gon row, written
  out because their support forms are transcendental-free. $\omega=2\pi f/32$ so that
  $f=1$ is one wave cycle per $32$ world units; $a=B/100$ for a bend slider $B$;
  $b=Ms/2\pi$ with $M=\max(1,\mathrm{round}(\lvert B\rvert/s))$ the number of
  spiral arms. The three concentric shapes are the support function of the
  generating polygon, which is the observation \S\ref{sec:walking} builds on.}
  \label{tab:catalog}
  \small
  \begin{tabular}{@{}llll@{}}
    \toprule
    Family & $\ph(p)$ & $\lvert\nabla\ph(p)\rvert$ & Notes \\
    \midrule
    Parallel lines
      & $\langle p,u\rangle$
      & $1$
      & $u=(\cos\alpha,\sin\alpha)$; pitch $s$ \\
    Concentric circles
      & $\lVert p\rVert$
      & $1$
      & \\
    Concentric squares
      & $\max(\lvert x\rvert,\lvert y\rvert)$
      & $1$
      & $L^\infty$; support function of the square \\
    Concentric triangles
      & $\max\!\left(x,\ \tfrac{\sqrt3}{2}\lvert y\rvert-\tfrac{x}{2}\right)$
      & $1$
      & support function of the triangle \\
    Concentric hexagons
      & $\max\!\left(\lvert x\rvert,\ \tfrac{\lvert x\rvert}{2}+\tfrac{\sqrt3}{2}\lvert y\rvert\right)$
      & $1$
      & support function of the hexagon \\
    Concentric $N$-gon
      & $\lVert p\rVert\cos\!\big(\mathrm{wrap}(\theta_p,\tfrac{2\pi}{N})\big)$
      & $1$
      & $\theta_p=\mathrm{atan2}(y,x)$ \\
    Wave
      & $x-A\sin(\omega y+\ph_0)$
      & $\sqrt{1+A^2\omega^2\cos^2(\omega y+\ph_0)}$
      & $A$ amplitude, $f$ frequency, $\ph_0$ phase \\
    Parabola
      & $y-a x^2$
      & $\sqrt{1+4a^2x^2}$
      & opens up; $a=0$ is horizontal parallels \\
    Hyperbola
      & $\sqrt{\lvert x^2-y^2\rvert}$
      & $\lVert p\rVert/\sqrt{\lvert x^2-y^2\rvert}$
      & rectangular; $n\ge1$, so no member through $0$ \\
    Spiral
      & $\lVert p\rVert-b\,\theta_p$
      & $\sqrt{1+b^2/\lVert p\rVert^2}$
      & Archimedean; $b=0$ is concentric circles \\
    \bottomrule
  \end{tabular}
\end{table*}

Three remarks on the table; each marks a bug we shipped.

\paragraph{The gradient column is not decoration.} For the first six rows it is
$1$ and can be dropped. For the last four it is not, and dropping it thins every
stroke by the same factor the family is steep by. We had shipped that bug in
three of the four; \S\ref{sec:eikonal} reports the measurement.

\paragraph{Some families exclude an index.} The hyperbola's $n=0$ member would be
the degenerate pair $x=\pm y$, and $\lvert\nabla\ph\rvert$ blows up there, so the
family starts at $n=1$. The spiral's $\ph$ has a branch cut wherever
$\mathrm{atan2}$ does; it closes only if the arm count $M$ is an integer, which is
why $M$ is rounded rather than passed through. Both are the kind of detail that a
raster pipeline hides in the resampling and a field pipeline exposes on the
first frame.

\paragraph{Three of the rows are support functions.} The square, triangle and
hexagon entries are all $\rad(q)=\max_k\langle q,n_k\rangle$ over the outward
facet normals, written out. So is the $N$-gon row, in its trigonometric form. This
is not a coincidence and it is not only a way to avoid an $\mathrm{atan2}$:
\S\ref{sec:walking} shows that the same object supplies the constants that make
the hard case tractable.

\subsection{Four families that are not phase functions}

\paragraph{The radial pencil.} $N$ undirected lines through the layer origin,
equally spaced over $\pi$. There is no $\ph$ here whose level sets are the members:
the index winds around the origin, so any candidate phase is circle-valued with
a winding-number obstruction at the core---a phase singularity in the sense of
Nye and Berry~\cite{Nye1974Dislocations}---and admits no global scalar lift.
What the family does have is a distance,
\begin{equation}
  d(p)=\max\!\Big(\lVert p\rVert\,\big\lvert\sin\mathrm{wrap}(\theta_p,\tfrac{\pi}{N})\big\rvert,
  \;\; \varphi-\lVert p\rVert\Big),
  \label{eq:radial}
\end{equation}
which is what the renderer needs (here and in Table~\ref{tab:catalog},
$\mathrm{wrap}(v,w)$ maps $v$ into $[-w/2,w/2)$; the formula is exact except
in the corner wedge where a line meets the opened disc, where it understates
the distance to the trimmed member by a bounded factor), and an index---the
angular sector $\lfloor \theta_p N/\pi\rfloor$---which is what the fringe
law needs. The second
term is the parameter the interface calls \emph{Start}: it opens a disc of radius
$\varphi$ around the origin, because $N$ lines crossing at a point is a black dot,
and opening the center is the standard fix. Because the sector index is piecewise
constant in $\theta_p$ rather than smooth, a pencil against a pencil does not
produce fringes in the sense of Theorem~\ref{thm:fringe}; it produces the rosettes
that the classical literature treats separately~\cite{Amidror2009Theory}.

\paragraph{Lattices.} Square, hexagonal and triangular lattices index their
members by \emph{two} integers, not one. The nearest lattice point is found by
rounding in the lattice basis---for the hexagonal and triangular cases, cube
rounding on $(u,v,-u-v)$, the $A_2$ closest-point rule of Conway and
Sloane~\cite{Conway1982Fast}, exact and branch-light---and the layer offers
two distances from the result: to the nearest vertex, and to the nearest edge. The
hexagonal edge distance is again a support form, $\min_k(\mathrm{apothem}-\langle
q,n_k\rangle)$ over the six sides, so hexagon edges are hexagon sides and not
three superposed line families; getting this wrong produces a plausible-looking
pattern with the wrong symmetry. Both distances take an anisotropic scale
$(s_x,s_y)$, applied by dividing by $\lVert(n_x/s_x,n_y/s_y)\rVert$ per facet so
that stretching a lattice keeps its strokes metrically correct.

The consequence for the theory is plain. Theorem~\ref{thm:fringe}
is about the difference of two \emph{scalar} index fields. A lattice's index is in
$\Z^2$, so the difference of two lattice indices is in $\Z^2$ as well, and the
theorem applies to it componentwise---once per generating line family. This is
the classical situation, where a $2$-D screen superposition has a fringe
system per pair of generating frequency vectors and the visible moir\'e is whichever
pair beats slowest~\cite{Amidror2009Theory, Amidror1998Fourier}. The field
formulation does not remove that bookkeeping. It localizes it: each component is a
scalar field we can evaluate, difference, and map.

\subsection{A family with no closed-form index}

Take any of the concentric rows of Table~\ref{tab:catalog} and give member $n$ its
own rigid motion---displace it by $n\delta$, turn it by $n\theta$. The picture is
still a family of nested shapes, the parameters are two sliders, and the result is
the geometry of much op art. But $\ph$ is gone: each member now lives in
its own frame, so no closed-form expression in $p$ has these curves as its level
sets. An index field still exists implicitly wherever consecutive members bracket
the point (\S\ref{sec:fringe}), which is why the fringe law still applies; what is
lost is the formula, so the nearest index must be searched for.

That is the subject of the next two sections, and it is where the field view stops
being a change of notation and starts costing something.
